% !TeX root = ../papers/nc-p3.tex
\section{八次モデルの有限境界における二次関係の極限}
\label{sec:rank-fifteen-octic-boundary}
\begin{theorem}\label{thm:rank-fifteen-octic-boundary}
$\xi^2+\xi+1=0$ の二水準と階数15の両表示について、標数零の通常窓で
$109$ における $K$ の剰余が $0,2,-2$ のいずれかとなるものは、
標準Koszul左右AS正則である。従って前節の有理 $t$ 剰余のうち、
未認証のものは $K=\pm1$ 上の
$5,22,37,53,56,72,87,104$ に限られる。
\end{theorem}
\begin{proof}
$F=t^8+14t^4+1$ とし、$Y^2=F$、$D=2t(t^2-1)$、$H=Y/D$ とおく。
$K=0$ では $t=\pm i$、$H=\pm i$ である。$K=k=\pm2$ では
$t=-k/2\pm\sqrt2$、$H=\pm\sqrt3$ である。これらの式は109での
全幾何剰余も尽くす。$Y$ は各点で単元なので、各標数零の境界点
$(\alpha,Y_0)$ の近傍では $u=t-\alpha$ が相対局所座標となる。
$\eta$ を原始24乗根として $i=\eta^6$、$j=\eta^2$、
$\sqrt2=\eta^3+\eta^{-3}$、$\sqrt3=\eta^2+\eta^{-2}$ を選ぶ。

\path{computations/verify_rank_fifteen_octic_boundary.py} は各位置・各文字の
関係行列から全最大小行列式を取り、一組のPl\"ucker座標全体の分母を一度に払う。
各成分は $\mathbb Z[i,\sqrt3][t,Y]$ の元となる。
$Y(0)=Y_0$ として $Y(u)^2=F(\alpha+u)$ から係数を順次求めると、
$2Y_0$ が単元であるため全係数が109の上で整である。各組において
全成分に共通な $u$ の最大冪を除く。実際の位数は各組で10未満であり、
これより低い係数の消滅は数体内で厳密に判定する。
最初のベクトルを109で還元すると、各組に非零な座標がある。
それをピボットに選んだGrassmann座標の行列は局所環上の直和因子を与え、
通常な持上げでは元の関係空間を与える。従ってこの行列の剰余は、
指定した円板のすべての持上げで共通である。

計算では $j$ の剰余 $8,41$ を取り、符号は二位置移動で尽くす。
$t\mapsto-1/t$、$Y\mapsto-Y/t^4$ の窓移動は $H$ を反転する。
$k=0$ では両 $t$ を取り、各 $H$ の一符号を調べる。
$k=\pm2$ では109上のFrobeniusが $\sqrt2$ を反転し $\sqrt3$ を保つので、
一つの $t,H$ からFrobeniusと窓移動で全四点を得る。
検算は各剰余の四位置で二次関係の階数6、経路次元35,56、
四次中心の文字別次元 $(3,2,0,0)$ を確認し、四つの中心元による商で
長さ13まで標準列を得る。通常持上げの長さ4,5の経路次元と
中心次元は定理\ref{thm:fixed-opposite-rank-drop}、\ref{thm:rank-fifteen-octic-center}
で同じである。従ってこれらの経路と中心の核は付値環上で基底変換と可換で、
中心切断を持ち上げられる。標準四切断判定を適用して結論を得る。
\end{proof}

\begin{proposition}
同じ通常表示の二次関係空間を基底を変えずに $K=\pm1$ へ延長すると、
その極限関係が生成する二次代数の長さ3,4,5の次元は、両水準の両表示と
全始点でそれぞれ $20,38,72$ となる。
\end{proposition}
\begin{proof}
各点は $Y=0$ にあり、$F'(\alpha)\ne0$ なので $u=Y$ を取る。
$U=A/K$、$V=B/K$ とすると
$\alpha=(V-1)/(U-1)$ で、$K=1$ では
$(U,V)=(\pm\sqrt3,\pm i)$、$K=-1$ では
$(U,V)=(\pm i,\pm\sqrt3)$ である。各Pl\"ucker座標を
$P(t)+YQ(t)$ と書く。$t-\alpha=u^2/F'(\alpha)+O(u^4)$ より、
$P,Q$ の最初の非零Taylor係数の次数をそれぞれ $a,b$ とすれば、
二項の位数は $2a,2b+1$ となる。零多項式は無限大の位数で扱う。
小さい方の係数を $F'(\alpha)^a$ または $F'(\alpha)^b$ で割ると
先頭係数を得る。\path{computations/verify_rank_fifteen_octic_path_jump.py}
は各組の最低位数の座標を取り、対応する関係行列で数体上の道を左に延ばす。
両符号の全組合せで長さ3,4,5の自由語を数えると主張の値になる。
他の $j$ は定数体の共役で従う。従って通常な一般ファイバーの
四次経路次元35を保つには、これら二つの境界で別途高次関係の飽和が要る。
\end{proof}
